%% ════════════════════════════════════════════════════════════════════════════
\section{Discussion}
\label{sec:discussion}

\subsection{Why Are Simple Baselines Hard to Beat in Gujarat?}
The empirical finding that complex machine learning algorithms fail to surpass simple classical baselines (with Persistence the simplest, and ARIMA(1,0,0) achieving lower MAE than or outperforming ML in Gandhinagar and Surat) across continental Gujarat is consistent with atmospheric persistence in these regions. In inland, semi-arid regions characterised by stable anticyclonic synoptic conditions during autumn and winter, daily boundary layer dynamics, nocturnal temperature inversions, and local anthropogenic emissions exhibit high day-to-day inertia. We hypothesise that because the atmospheric system changes gradually, yesterday's 24-hour mean AQI represents a near-sufficient statistic for tomorrow's mean state under unperturbed conditions. Machine learning models trained on purely lagged surface concentrations cannot anticipate external synoptic shifts without numerical meteorological forcing.

\subsection{The Coastal Exception: Dynamics of Chennai}
In contrast, Chennai's maritime boundary layer is governed by diurnal mesoscale land-sea breeze circulations and maritime moisture advection. These dynamics may reduce day-to-day AQI inertia, creating structured periodic patterns that cyclical calendar harmonics and autoregressive AQI patterns can exploit (noting that Chennai's 2022--2023 test window has imputed pollutant lags). This may explain why machine learning models achieved nominally significant $+11.5\%$ skill improvements in Chennai while failing to beat persistence in Ahmedabad and Gandhinagar; we treat this as a working hypothesis pending residual analysis with meteorological covariates.

\subsection{Overfitting to Industrial Emissions: The Case of Vapi}
The statistically significant underperformance of Random Forest and XGBoost in Vapi ($p \leq 0.003$ relative to Persistence) highlights a critical vulnerability in applied environmental machine learning. Vapi is an intensive chemical manufacturing hub subject to sporadic, unannounced industrial batch releases. These intermittent emissions produce heavy-tailed spikes in criteria pollutants. We hypothesise that nonlinear tree ensembles may branch heavily on these high-magnitude anomalies during training, learning spurious split boundaries. When ambient conditions revert to typical baselines, these models project lingering elevations, producing substantial over-prediction errors. Simple persistence, which tracks only the immediate preceding state, avoids this structural over-prediction.

\subsection{Recommendations for Environmental Data Science}
Based on our findings, we propose four methodological standards for atmospheric machine learning studies:
\begin{enumerate}
  \item \textbf{Persistence and AR(1) Benchmarking:} Every urban AQI forecasting study should include the naive Persistence baseline ($\hat{y}_{t+h} = y_t$) and a simple AR(1) benchmark, reporting relative forecast skill scores rather than raw accuracy metrics alone.
  \item \textbf{Avoid Random Splitting:} Random train-test partitioning should not be used on environmental time series. Evaluations should enforce temporal separation via rolling-origin or blocked prequential protocols.
  \item \textbf{Exclusion of Concurrent Pollutants from Forecasting Benchmarks:} Predictor matrices should exclude same-day criteria pollutant concentrations unless the task is explicitly framed as algebraic reconstruction or nowcasting.
  \item \textbf{Formal Statistical Significance Testing:} Authors should conduct formal hypothesis tests (e.g., Diebold-Mariano with HAC variance estimators) to substantiate claims of algorithmic superiority.
\end{enumerate}

\subsection{Limitations}
\label{sec:limitations}
Several limitations of the present study should be noted:
\begin{itemize}
  \item \textbf{Station Network Heterogeneity:} City-level AQI values represent spatial averages across active municipal CAAQMS stations. Station additions or maintenance outages over the 2015--2023 timeframe may introduce minor unmodelled level shifts in the aggregate series.
  \item \textbf{Exclusion of Numerical Weather Prediction (NWP) Forcing:} To isolate the autoregressive information content of surface monitoring records, this benchmark excluded numerical weather forecasts (wind vector, mixing height, relative humidity). Incorporating high-resolution operational NWP grids represents an essential avenue for future research.
  \item \textbf{Fixed Hyperparameter Regimes:} Primary benchmark hyperparameters were held constant across cities to prevent post-hoc optimization leakage. While a 3-fold \texttt{TimeSeriesSplit} tuned check on Ahmedabad yielded no significant improvement over Persistence ($\text{MAE} = 19.99\text{--}21.03$, all $p \ge 0.30$), full nested rolling-origin tuning across all cities could yield marginal refinements.
  \item \textbf{Pollutant Coverage and Regime Shifts:} Kaggle pollutant records end in 2020, so training folds in Ahmedabad and Chennai contain observed 2015--2020 pollutant lags whereas their 2022--2023 Setup~C test inputs are imputed. A 21-column AQI-only ablation confirms that Ahmedabad ML still fails to beat Persistence ($\text{MAE} = 21.17\text{--}21.33$, $p \ge 0.168$), while in Chennai Random Forest drops from $15.84$ ($p = 0.002$) to $16.84$ ($p = 0.087$). Similarly, Setup~B's test window overlaps the March--May 2020 COVID-19 lockdown regime shift, which accounts for most of its apparent $R^2$ collapse.
  \item \textbf{Geographical Scope:} Our empirical scope centered on western India (Gujarat) and a southern maritime anchor (Chennai). Extending this audit to the Indo-Gangetic Plain, where intense seasonal biomass burning and winter fog inversions prevail, constitutes an important next step.
\end{itemize}

